\scanpage{302}
\mspage{218}

tel qu'on ait

\[
\varphi\xleftarrow{\sim}\varinjlim_{i\in I}\Phi(i),\tag*{(4.9.5)}
\]

où ce qui revient au \unsure{même}, tel qu'on ait

\[
\varphi(j)\xleftarrow{\sim}\varinjlim_{i\in I}\Phi(i,j),\quad\text{pour }j\in J\tag*{(4.9.6)}
\]

où on a posé

\[
\Phi(i,j)=\Phi(i)(j)
\]

interprétant donc $\Phi$ comme un foncteur

\[
\Phi:I\times J\longrightarrow C\ ;\tag*{(4.9.7)}
\]

la loi fonctorielle de $\varphi(j)$ donnée par (4.9.6) \unsure{étant|est} évidente. En d'autres termes, les $\varphi$ dans l'image essentielle de (4.9.3) sont ceux qui admettent une ``\uline{représentation indicielle}'' par un foncteur $\Phi$ sous forme (4.9.4), ou sous forme (4.9.7), avec \uline{de plus} $I$ grd devant $\pi_{\unsure{0|1}}$. Je rappelle à ce sujet la

\uline{Proposition} 4.9.2. Sous les conditions précédentes sur $C,\pi_0,J$ (notamment $\operatorname{card}\operatorname{réd}J\leq\pi_0$), si $\varphi:J\longrightarrow M\overset{\mathrm{\unsure{def|déf}}}{=}\operatorname{Ind}_{\pi_0}(C)$, pour que $\varphi$ appartienne à l'image essentielle du foncteur pl. fid. (4.9.3), il f. et s. qu'elle satisfasse les deux conditions suivantes :
% archive sheet number: 301
